\chapter{双元法与mod法}
\section{mod法}
有理分式分解 $\frac{P(x)}{A(x)B(x)}=\frac{U(x)}{A(x)} + \frac{V(x)}{B(x)}$。通分，等价于求解：
$$ U(x)B(x) + V(x)A(x) = P(x) $$
在等式两边对 $A(x)$ 取模（也就是令 $A(x)=0$），那么 $V(x)A(x)$ 这一项消失，式子变成：
$$ U(x)B(x) \equiv P(x) \pmod{A(x)} \implies U(x) \equiv P(x) \cdot \color{white}{B^{-1}(x)} \pmod{A(x)}$$
问题转化为求 $B(x)$ 在模 $A(x)$ 下的“乘法逆元” $B^{-1}(x)$（也就是多项式的倒数）。
只要 $A(x)$ 和 $B(x)$ 没有公因式，就存在多项式 $u(x)$ 和 $v(x)$ 使 $u(x)A(x) + v(x)B(x) = 1$，因此逆元总能从裴蜀等式里构造出来。遇到重因式 $(x^2+1)^2$ 时，可以把分子连续对 $(x^2+1)$ 取余，逐步拆开；若分母含混合重因式 $A^m B^n$，则可从一次情形的裴蜀等式出发，再把等式两边取 $(m+n-1)$ 次方展开，这里用到的就是这种二项式升幂法。

\begin{example}[混合分解 $\frac{x^3+2}{(x-1)^2(x^2+1)}$]
    分解 $\displaystyle F(x) = \frac{x^3+2}{(x-1)^2(x^2+1)}$
\end{example}
\begin{solution}
原分母含重根线性因式 $(x-1)^2$ 与二次因式 $(x^2+1)$。先处理 $(x-1)^2$。设分解形式中包含 $\frac{A}{x-1} + \frac{B}{(x-1)^2}$，令剩余部分函数为 $G(x) = \frac{x^3+2}{x^2+1}$。
二阶极点的最高阶系数来自 \((x-1)^2F(x)\) 在 \(x=1\) 的值，因此代入 \(x=1\) 得
$$ B = G(1) = \frac{1^3+2}{1^2+1} = \frac{3}{2}. $$
再对 $G(x)$ 求一阶导数并代入 $x=1$，
\begin{align*}
G'(x) &= \frac{3x^2(x^2+1) - (x^3+2)(2x)}{(x^2+1)^2} ,~~A = G'(1) = \frac{3(2) - (3)(2)}{2^2} = 0
\end{align*}
由于导数值为 \(0\)，一阶项消失，重根部分目前只保留 \(\frac{3/2}{(x-1)^2}\)。从原式中减去这项，记剩余部分为 $H(x)$，则
\begin{align*}
H(x) = \frac{x^3+2}{(x-1)^2(x^2+1)} - \frac{3/2}{(x-1)^2} = \frac{x^3+2 - \frac{3}{2}(x^2+1)}{(x-1)^2(x^2+1)} = \frac{x^3 - \frac{3}{2}x^2 + \frac{1}{2}}{(x-1)^2(x^2+1)}.
\end{align*}
这时分子必含 $(x-1)^2$ 因式，因而
$$ x^3 - \frac{3}{2}x^2 + \frac{1}{2} = (x-1)^2 \left(x + \frac{1}{2}\right). $$
代回并约去 $(x-1)^2$，得到
$$ H(x) = \frac{(x-1)^2 \left(x + \frac{1}{2}\right)}{(x-1)^2(x^2+1)} = \frac{x + \frac{1}{2}}{x^2+1}. $$
于是
$$ F(x) = \frac{3}{2(x-1)^2} + \frac{2x+1}{2(x^2+1)}. $$
（注：若 $H(x)$ 的分母中仍包含多个高次不可约多项式，可以继续对相应因式取余；这样每一步只处理低次多项式，避免复数根和大方程组。）
\end{solution}

\begin{example}[分解$\frac{x+2}{(x^2+2x+2)(x^2+1)}$]
分解$\dfrac{x+2}{(x^2+2x+2)(x^2+1)}$
\end{example}
\begin{solution}
设大分母 $A = x^2+2x+2$，小分母 $B = x^2+1$。目标是将它们线性组合成 $1$。
先用 $A$ 除以 $B$，得到 $A \equiv 2x+1 \pmod{x^2+1}$；再用 $B=x^2+1$ 除以余式 $2x+1$，可得商 $\frac{2x-1}{4}$，余数 $\frac{5}{4}$，于是
\[
B - (2x+1)\left(\frac{2x-1}{4}\right) = \frac{5}{4}.
\]
再把余式 $2x+1$ 用 $A-B$ 替换，就有
\[
\frac{5}{4} = B - (A - B)\left(\frac{2x-1}{4}\right)\Leftrightarrow 5 = B\left(2x+3\right) - A\left(2x-1\right).
\]
原式因此可写成
\begin{align*}
\frac{x+2}{AB}=&\frac{(x+2)5}{5AB}=\frac{(x+2)\left( B\left(2x+3\right) - A\left(2x-1\right)\right)}{5AB}\\
=&\frac{(x+2)(2x+3)}{5(x^2+2x+2)} - \frac{(x+2)(2x-1)}{5(x^2+1)}=\frac{2x^2+7x+6}{5(x^2+2x+2)} - \frac{2x^2+3x-2}{5(x^2+1)}\\
=&\frac{2x^2+7x+6}{5(x^2+2x+2)}-\frac25 - \left(\frac{2x^2+3x-2}{5(x^2+1)}-\frac25\right)\\
=&\frac{4-3x}{5(x^2+1)}+ \frac{3x+2}{5(x^2+2x+2)}.
\end{align*}
\end{solution}

\begin{example}[提取 $\frac{x^3+3x}{(x^2+1)^2}$]
提取 $\dfrac{x^3+3x}{(x^2+1)^2}$
\end{example}
\begin{solution}
先把 \(x^3+3x\) 写成 \(x(x^2+1)+2x\)，再除以 \((x^2+1)^2\)，得到
\begin{align*}
\frac{x(x^2+1) + 2x}{(x^2+1)^2} = {\frac{x}{x^2+1} + \frac{2x}{(x^2+1)^2}} 
\end{align*}
\end{solution}

\begin{example}[分解$\frac{1}{(x-1)^2(x^2+1)^2}$]
分解$\dfrac{1}{(x-1)^2(x^2+1)^2}$
\end{example}
\begin{solution}
设 $A = x-1, B = x^2+1$。
先算得 $B$ 除以 $A$ 的商为 $x+1$，余式为 $2$，即
$$ B - A(x+1) = (x^2+1) - (x^2-1) = 2. $$
分母是 $A^2B^2$，指数之和减一为 $2+2-1=3$，因此把上式两边立方：
$$ \left(B - (x+1)A\right)^3=8. $$
展开后有
$$ B^3-3(x+1)AB^2+3(x+1)^2A^2B-(x+1)^3A^3=8, $$
再按 $B^2$ 和 $A^2$ 分组，
$$ 8 = B^2 \underbrace{\left[ B -3(x+1)A \right]}_{\text{这是} A^2 \text{的分子} U(x)} + A^2 \underbrace{\left[ {3(x+1)^2}B - (x+1)^3A \right]}_{\text{这是} B^2\text{ 的分子} V(x)}. $$
代入 $A = x-1, B = x^2+1$ 展开括号，得到
$$ 8=B^2(4-2x^2)+A^2(2x^4+4x^3+6x^2+8x+4). $$
继续做竖式除法，
\begin{align*}
U(x) \div (x-1)^2 &\implies 4-2x^2 =-2(x-1)^2 + 6-4x\\
V(x) \div (x^2+1)^2 &\implies 2x^4+4x^3+6x^2+8x+4= 2(x^2+1)^2 +4x^3+2x^2+8x+2,
\end{align*}
把两次降次得到的整式部分分别写出后，常数项一正一负相抵，剩余部分为
$$ \frac{3-2x}{4(x-1)^2} + \frac{2x^3+x^2+4x+1}{4(x^2+1)^2}. $$
\end{solution}



\begin{example}[分解 $\frac{2x^2 + 3}{(x-1)(x^2+x+1)}$]
    分解 $\displaystyle F(x) = \frac{2x^2 + 3}{(x-1)(x^2+x+1)}$
\end{example}
\begin{solution}
设二次因式 $A = x^2+x+1$，一次因式 $B = x-1$。
用带余除法可得 $A = B(x+2) + 3$，于是
$$ 1 = \frac{1}{3}A - \frac{x+2}{3}B. $$
把它乘到原分式分子 $P(x)=2x^2+3$ 上，再除以总分母 $AB$，便有
\begin{align*}
\frac{P(x)}{AB} &= \frac{P(x) \left(\frac{1}{3}A - \frac{x+2}{3}B\right)}{AB} = \frac{2x^2+3}{3B} - \frac{(2x^2+3)(x+2)}{3A}.
\end{align*}
接着分别把两项分子降次：
\begin{align*}
\frac{2x^2+3}{3(x-1)} &= \frac{2(x^2-1)+5}{3(x-1)} = \frac{2(x+1)}{3} + \frac{5}{3(x-1)} \\
\frac{(2x^2+3)(x+2)}{3(x^2+x+1)} &= \frac{2x^3+4x^2+3x+6}{3(x^2+x+1)} = \frac{2x(x^2+x+1) + 2(x^2+x+1) - x + 4}{3(x^2+x+1)} \\
&= \frac{2x+2}{3} + \frac{-x+4}{3(x^2+x+1)}.
\end{align*}
两式相减时，多项式部分自然抵消，最终得到
\begin{align*}
F(x) &= \left(\frac{2x+2}{3} + \frac{5}{3(x-1)}\right) - \left(\frac{2x+2}{3} + \frac{-x+4}{3(x^2+x+1)}\right) \\
&= \frac{5}{3(x-1)} + \frac{x-4}{3(x^2+x+1)}.
\end{align*}
\end{solution}

\begin{example}[分解 $\frac{x^5}{(x^2+1)^3}$]
    分解 $\displaystyle F(x) = \frac{x^5}{(x^2+1)^3}$
\end{example}
\begin{solution}
令 $t = x^2+1$，则 $x^2 = t-1$。于是
$$ x^5 = x \cdot (x^2)^2 = x(t-1)^2 = x(t^2 - 2t + 1) = x \cdot t^2 - 2x \cdot t + x. $$
把 $t$ 换回 $x^2+1$ 并代回原式，连同分母 $(x^2+1)^3$ 逐项约分，便得
\begin{align*}
F(x) &= \frac{x(x^2+1)^2 - 2x(x^2+1) + x}{(x^2+1)^3} \\
&= \frac{x}{x^2+1} - \frac{2x}{(x^2+1)^2} + \frac{x}{(x^2+1)^3}.
\end{align*}
\end{solution}

\begin{example}[分解 $\frac{x}{(x-1)^2(x+1)^2}$]
    分解 $\displaystyle F(x) = \frac{x}{(x-1)^2(x+1)^2}$
\end{example}
\begin{solution}
设 $A = x+1$，$B = x-1$。
由 $A - B = 2$ 可得
$$ 1 = \frac{A - B}{2}. $$
原分母含 $A^2B^2$，总指数为 $4$，于是把上式两边升到 $3$ 次方：
$$ 1 = \left(\frac{A - B}{2}\right)^3 = \frac{1}{8}(A^3 - 3A^2B + 3AB^2 - B^3). $$
按 $A^2$ 和 $B^2$ 提取公因式后，
$$ 1 = A^2 \left( \frac{A - 3B}{8} \right) + B^2 \left( \frac{3A - B}{8} \right). $$
再代入 $A=x+1, B=x-1$ 化简，得到
$$ 1 = A^2 \left( \frac{2-x}{4} \right) + B^2 \left( \frac{x+2}{4} \right). $$
把它代回原分式分子并与分母 $A^2B^2$ 约分，
$$ F(x) = \frac{x \cdot 1}{A^2 B^2} = \frac{x(2-x)}{4 B^2} + \frac{x(x+2)}{4 A^2} = \frac{2x-x^2}{4(x-1)^2} + \frac{x^2+2x}{4(x+1)^2}. $$
两项分别降次为
\begin{align*}
\frac{2x-x^2}{4(x-1)^2} &= \frac{-(x^2-2x+1)+1}{4(x-1)^2} = -\frac{1}{4} + \frac{1}{4(x-1)^2} \\
\frac{x^2+2x}{4(x+1)^2} &= \frac{(x^2+2x+1)-1}{4(x+1)^2} = \frac{1}{4} - \frac{1}{4(x+1)^2},
\end{align*}
相加后常数项抵消，因此
$$ F(x) = \frac{1}{4(x-1)^2} - \frac{1}{4(x+1)^2}. $$
\end{solution}

\begin{example}[分解 $\frac{x^2}{(x^2+1)^2(x^2+x+1)}$]
    分解 $\displaystyle F(x) = \frac{x^2}{(x^2+1)^2(x^2+x+1)}$
\end{example}
\begin{solution}
设 $A = x^2+x+1$，$B = x^2+1$。原式分母结构为 $A B^2$。
由 $x^2=B-1$ 可先写成
$$ \frac{x^2}{A B^2} = \frac{B - 1}{A B^2} = \frac{1}{AB} - \frac{1}{A B^2}, $$
于是原式先拆成 \(\frac{1}{AB}\) 与 \(\frac{1}{AB^2}\) 两块。由 $A-B=x$，再利用带余除法 $B \div x$ 得 $B = x \cdot x + 1$，把 $x=A-B$ 代入余式关系便有
$$ 1 = B - x(A-B) = (x+1)B - xA, $$
从而
$$ \frac{1}{AB} = \frac{(x+1)B - xA}{AB} = \frac{x+1}{A} - \frac{x}{B}. $$
把它两边同除以 $B$，可得
$$ \frac{1}{AB^2} = \frac{x+1}{AB} - \frac{x}{B^2}. $$
再将恒等式 $1 = (x+1)B - xA$ 两端同乘 $(x+1)$，
$$ x+1 = (x+1)^2 B - x(x+1)A, $$
除以 $AB$ 后得到
$$ \frac{x+1}{AB} = \frac{(x+1)^2}{A} - \frac{x^2+x}{B}, $$
于是
$$ \frac{1}{A B^2} = \frac{(x+1)^2}{A} - \frac{x^2+x}{B} - \frac{x}{B^2}. $$
再把分子降次：
\begin{align*}
\frac{(x+1)^2}{A} &= \frac{x^2+2x+1}{x^2+x+1} = \frac{(x^2+x+1)+x}{x^2+x+1} = 1 + \frac{x}{A} \\
\frac{x^2+x}{B} &= \frac{x^2+x}{x^2+1} = \frac{(x^2+1)+x-1}{x^2+1} = 1 + \frac{x-1}{B},
\end{align*}
代回便是
$$ \frac{1}{A B^2} = \left(1 + \frac{x}{A}\right) - \left(1 + \frac{x-1}{B}\right) - \frac{x}{B^2} = \frac{x}{A} - \frac{x-1}{B} - \frac{x}{B^2}. $$
最后相减合并：
\begin{align*}
F(x) &= \frac{1}{AB} - \frac{1}{AB^2} = \left(\frac{x+1}{A} - \frac{x}{B}\right) - \left(\frac{x}{A} - \frac{x-1}{B} - \frac{x}{B^2}\right) \\
&= \frac{x+1-x}{A} + \frac{-x+(x-1)}{B} + \frac{x}{B^2} \\
&= \frac{1}{x^2+x+1} - \frac{1}{x^2+1} + \frac{x}{(x^2+1)^2}.
\end{align*}
\end{solution}

\begin{example}[分解 $F(x) = \frac{1}{x(x-1)(x^2+1)}$]
    分解 $\displaystyle F(x) = \frac{1}{x(x-1)(x^2+1)}$
\end{example}
\begin{solution}
先拆出因式 $x$。设 $A = x$，其余部分 $M = (x-1)(x^2+1) = x^3 - x^2 + x - 1$。由带余除法知 $M \div A$ 的商为 $x^2-x+1$，余数为 $-1$，即
$$ M = A(x^2-x+1) - 1. $$
于是
$$ 1 = A(x^2-x+1) - M, $$
代入原分式并裂项可得
$$ \frac{1}{AM} = \frac{A(x^2-x+1) - M}{AM} = \frac{x^2-x+1}{M} - \frac{1}{A} = \frac{x^2-x+1}{(x-1)(x^2+1)} - \frac{1}{x}. $$
再分解新分式 $\frac{x^2-x+1}{(x-1)(x^2+1)}$。设 $B = x^2+1$，$C = x-1$，则 $B \div C$ 的商为 $x+1$，余数为 $2$，也就是
$$ B = C(x+1) + 2. $$
因此可构造
$$ 1 = \frac{1}{2}B - \frac{x+1}{2}C, $$
把它乘入分子 $P(x) = x^2-x+1$ 并除以 $BC$，有
\begin{align*}
\frac{P(x)}{BC} &= \frac{P(x)\left(\frac{1}{2}B - \frac{x+1}{2}C\right)}{BC} = \frac{x^2-x+1}{2(x-1)} - \frac{(x^2-x+1)(x+1)}{2(x^2+1)}.
\end{align*}
两部分分别降次：
\begin{align*}
\frac{x^2-x+1}{2(x-1)} &= \frac{x(x-1)+1}{2(x-1)} = \frac{x}{2} + \frac{1}{2(x-1)} \\
\frac{x^3+1}{2(x^2+1)} &= \frac{x(x^2+1)-x+1}{2(x^2+1)} = \frac{x}{2} + \frac{1-x}{2(x^2+1)}.
\end{align*}
相减后一次式自然抵消，
$$ \left(\frac{x}{2} + \frac{1}{2(x-1)}\right) - \left(\frac{x}{2} + \frac{1-x}{2(x^2+1)}\right) = \frac{1}{2(x-1)} + \frac{x-1}{2(x^2+1)}, $$
故原分式分解为
$$ F(x) = -\frac{1}{x} + \frac{1}{2(x-1)} + \frac{x-1}{2(x^2+1)}. $$
\end{solution}

\begin{example}[分解 $\frac{x^3}{(x-1)(x+1)(x^2+x+1)}$]
    分解 $\displaystyle F(x) = \frac{x^3}{(x-1)(x+1)(x^2+x+1)}$
\end{example}
\begin{solution}
对单根项，先把对应因子乘回原分式，再令该因子为零取值，这样得到该简单极点的系数。令 $x=1$，掩去分母中的 $(x-1)$，代入剩余部分得
$$ \frac{1^3}{(1+1)(1+1+1)} = \frac{1}{6}, $$
所以含 $\frac{1/6}{x-1}$；令 $x=-1$，掩去分母中的 $(x+1)$，代入剩余部分得
$$ \frac{(-1)^3}{(-1-1)(1-1+1)} = \frac{1}{2}, $$
所以含 $\frac{1/2}{x+1}$。先把这两项合并：
$$ \frac{1}{6(x-1)} + \frac{1}{2(x+1)} = \frac{1(x+1) + 3(x-1)}{6(x^2-1)} = \frac{4x-2}{6(x^2-1)} = \frac{2x-1}{3(x^2-1)}. $$
从原式中减去这部分，剩余项为
$$ \text{剩余项} = \frac{x^3}{(x^2-1)(x^2+x+1)} - \frac{2x-1}{3(x^2-1)}. $$
通分相减后，分子变成
$$ 3x^3 - (2x-1)(x^2+x+1) = 3x^3 - (2x^3+x^2+x-1) = x^3 - x^2 - x + 1, $$
而它恰好可以写成
$$ x^3 - x^2 - x + 1 = x^2(x-1) - (x-1) = (x^2-1)(x-1). $$
代回并约去 $(x^2-1)$，得到
$$ \text{剩余项} = \frac{(x^2-1)(x-1)}{3(x^2-1)(x^2+x+1)} = \frac{x-1}{3(x^2+x+1)}. $$
因此最终结果为
$$ F(x) = \frac{1}{6(x-1)} + \frac{1}{2(x+1)} + \frac{x-1}{3(x^2+x+1)}. $$
\end{solution}


\section{双元法}
本章省略了常数$C$。
\subsection{双元微分公式}
\begin{theorem}[微分公式]
实圆：$x\dd x+y\dd y=0$；虚圆：$x\dd x=y\dd y$。参数微元方面，实圆满足 $\frac{\dd x}{y}=-\dd\theta, \quad \frac{\dd y}{x}=\dd\theta$，虚圆满足 $\frac{\dd x}{y}=\dd\phi,\quad \frac{\dd y}{x}=\dd\phi$。面积 $2$ 倍的微元可写成 $y\dd x-x\dd y=(x^2\pm y^2)\dd\theta$，这里要注意顺序和微分的顺序一致；比值微元则是 $\dd\left(\frac{x}{y}\right)=\frac{y\dd x-x\dd y}{y^2}$。
\end{theorem}
实圆关系 $x\dd x+y\dd y=0$ 的几何意义是，瞬时位移向量 $(\dd x, \dd y)$ 必须垂直于位置向量 $(x,y)$，二者内积为零；虚圆关系 $x\dd x=y\dd y$ 则等价于 $\frac{\dd x}{x} = \frac{\dd y}{y}$，也就是 $x$ 与 $y$ 的相对变化率必须相等。

参数微元也可从几何关系得到。对实圆，$\dd\theta$ 是位置向量绕原点逆时针转过的有向小角，此时指针尖端的位移方向为 $\vec{r}$ 逆时针旋转 $90^\circ$ 后得到的方向，坐标为 $(-y,x)$，长度为弧长 $a\dd\theta$，所以
\[(\dd x, \dd y)=\dd\theta(-y,x)\]
对虚圆，$\dd\phi$ 是位置向量沿等相对变化率运动时的参数微元。由 \(x\dd x=y\dd y\) 可取 \((\dd x,\dd y)\) 与 \((y,x)\) 同向，长度由参数微元吸收进 \(\dd\phi\)，所以
\[(\dd x, \dd y)=\dd\theta(y,x)\]

面积 $2$ 倍的微元来自扇形面积两倍等于半径平方乘以角度，方向则要随着微分顺序一起看。因此
\[
x\dd y-y\dd x=(x^2\pm y^2)\dd\theta
\]
表示逆向面积，对应逆时针角度微元 $\dd\theta$，也就是实圆关系中的 $\frac{\dd y}{x}$，以及虚圆关系中的 $\frac{\dd y}{x},\frac{\dd x}y$；而
\[
y\dd x-x\dd y=(x^2\pm y^2)(-\dd\theta)
\]
表示顺向面积，对应顺时针角度微元 $-\dd\theta$，也就是实圆关系中的 $\frac{\dd x}{y}$。实际计算时，可以先补出分母里的 \(x^2\pm y^2\)，再观察分子与分母的对应关系：$\frac{x\dd y-y\dd x}{x^2+y^2}$ 对应 $\frac{\dd y}{x},-\frac{\dd x}{y}$ 的合分比，$\frac{y\dd x-x\dd y}{x^2+y^2}$ 对应 $\frac{\dd x}{y},-\frac{\dd y}{x}$ 的合分比，$\frac{x\dd y-y\dd x}{x^2-y^2}$ 对应 $\frac{\dd y}{x},\frac{\dd x}{y}$ 的合分比，$\frac{y\dd x-x\dd y}{x^2-y^2}$ 则对应 $-\frac{\dd x}{y},-\frac{\dd y}{x}$ 的合分比。

比值微元 $\dd\left(\frac{x}{y}\right)=\frac{y\dd x-x\dd y}{y^2}$ 经常反过来用作积分降次，也可以用相似三角形来理解。取参考线 $y=1$，将原点与动点 $P(x,y)$ 连成一条射线，交参考线于 $P'$，由相似三角形可知 $P'$ 的坐标正是 $\left(\frac{x}{y}, 1\right)$，因此 $\frac{x}{y}$ 就是辅助点的横坐标。当 $P'$ 在直线上发生水平微移 $\dd\left(\frac{x}{y}\right)$ 时，它与原点围成的小三角形按 $2$ 倍有向面积计算满足
\[
\frac{x}{y}\cdot 0 - 1\cdot\dd\left(\frac{x}{y}\right) = -\dd\left(\frac{x}{y}\right).
\]
另一方面，实际动点 $P(x,y)$ 与辅助点 $P'$ 在同一条射线上，长度比为 $y:1$，所以对应面积按相似比的平方放大为 $y^2$ 倍，于是
\[ x\dd y - y\dd x = y^2 \left[ -\dd\left(\frac{x}{y}\right) \right]. \]
把负号移入括号整理，便得到 $\dd\left(\frac{x}{y}\right)=\frac{y\dd x-x\dd y}{y^2}$。

\subsection{双元积分公式}
\begin{theorem}[基本积分公式]
实圆：$\int\frac{\dd x}{y}=\arctan\frac{x}{y}$；虚圆：$\int \frac{\dd x}{y}=\ln(x+y)$。
\end{theorem}
对于虚圆 \(x^2 - y^2 = 1\)（或一般常数），参数方程可写为 \(x = \cosh \phi\)，\(y = \sinh \phi\)，那么：
\[
x + y = \cosh \phi + \sinh \phi = e^\phi,\quad\phi = \ln(x + y).
\]
而微分 \(\frac{\dd x}{y} = \dd\phi\)（因为 \(\dd x = \sinh \phi \, \dd\phi = y \, \dd\phi\)），因此
\[
\int \frac{dx}{y} = \phi + C = \ln(x + y) + C.
\]
在双曲线上，\(\phi\) 可视为双曲角参数，而 \(x+y\) 是相应的指数坐标，故 \(\ln(x+y)\) 正是该参数的表达式。

对于实圆 \(x^2 + y^2 = 1\)，可取参数 \(x = \cos\theta\)，\(y = \sin\theta\)。此时积分：
\[\frac{\dd x}{y} = -\dd \theta\Rightarrow\int \frac{\dd x}{y} = -\theta + C\]
另一方面，比值 \(\frac{x}{y} = \cot\theta\)，所以
\[\arctan\frac{x}{y} = \arctan(\cot\theta) = \frac{\pi}{2} - \theta \quad (\text{在适当范围内})\]
因此\[\int \frac{dx}{y} = \arctan\frac{x}{y} + \left(C - \frac{\pi}{2}\right).\]
即两者相差常数，所以不定积分公式成立。

\begin{theorem}[三次倒数型]
\vspace{-15pt}
\begin{align*}
\int \frac{\dd x}{y^3} &=\int\frac{1}{y^2}\frac{\dd x}{y}=\int\frac{1}{y^2}\frac{y\dd x-x\dd y}{y^2\pm x^2} =\frac{1}{y^2\pm x^2}\int\frac{y\dd x-x\dd y}{y^2}=\frac{1}{y^2\pm x^2}\int\dd\left(\frac{x}{y}\right)\\
\int \frac{\dd x}{xy^2}&=\int\frac1{xy}\frac{\dd x}{y}=\int\frac1{xy}\frac{y\dd x-x\dd y}{y^2\pm x^2}=\frac{1}{y^2\pm x^2}\int\left(\frac{\dd x}{x}-\frac{\dd y}{y}\right)\\
\int \frac{\dd x}{x^2y}&=\int\frac{1}{x^2}\frac{\dd x}{y}=\int\frac{1}{x^2}\frac{y\dd x-x\dd y}{y^2\pm x^2}=\frac{-1}{y^2\pm x^2}\int\dd\left(\frac{y}{x}\right)
\end{align*}
\end{theorem}
遇到的式子若缺少匹配的面积微元或比值微元，就先补出 \(y\dd x-x\dd y\) 这类结构，再借助微分公式完成转化。

\begin{theorem}[线性微元型]
\vspace{-15pt}
\begin{align*}
\int x\dd y&=\int\frac12\left(x\dd y+y\dd x+x\dd y-y\dd x\right)=\frac12\int\dd(xy)+\frac{x^2\pm y^2}2\int\dd\frac{x}{y}\\
\int y\dd x&=\int\frac12\left(y\dd x+x\dd y+y\dd x-x\dd y\right)=\frac12\int\dd(xy)+\frac{x^2\pm y^2}2\int\dd\frac{y}{x}
\end{align*}
\end{theorem}
拆分的目标是把一组整理成 \(\dd(xy)\)，另一组整理成面积微元，再分别积分。

\begin{theorem}[混合微元型]
\vspace{-15pt}
\begin{align*}
\int \frac{\dd x}{xy}&=\frac{1}{C}\int\frac{x^2\pm y^2}{xy}\dd x=\frac1{C}\left(\int\frac{x\dd x}{y}\pm \int\frac{y\dd x}{x}\right)\\
\int\frac{x\dd y}{y}&=\int\frac{xy\dd y}{y^2}
\end{align*}
之后把 $y\dd y$ 和分母的 $y^2$ 换掉，化成前面的一元形式。
\end{theorem}

\begin{theorem}[负偶次幂型]
负偶数次幂，用“常数代换+微分代换+分部积分”进行降次。
\end{theorem}
例如
\begin{align*}
\int\dfrac{\dd x}{y^4}&=\frac{1}{C}\int\frac{x^2\pm y^2}{y^4}\dd x=\frac{1}{C}\left(\int\frac{x^2\dd x}{y^4}\pm\int\frac{\dd x}{y^2}\right)\\
&=\begin{cases}\frac{1}{C}\left(\int\frac{xy\dd y}{y^4}\pm\int\frac{\dd x}{y^2}\right)=\frac{1}{C}\left(\int\frac{x\dd y}{y^3}\pm\int\frac{\dd x}{y^2}\right)\\\frac{1}{C}\left(\int\frac{-xy\dd y}{y^4}\pm\int\frac{\dd x}{y^2}\right)=\frac{1}{C}\left(\int\frac{-x\dd y}{y^3}\pm\int\frac{\dd x}{y^2}\right)\end{cases}
\end{align*}
对 $\displaystyle\int\frac{x\dd y}{y^3}$ 用分部积分
\begin{align*}
    \int\frac{x\dd y}{y^3}&=\int x\dd\left(\frac1{-2}\frac{1}{y^2}\right)=-\frac12\int x\dd\left(\frac{1}{y^2}\right)\\
    &=-\frac12\frac{x}{y^2}+\frac12\int\frac{\dd x}{y^2}\\
\end{align*}
更高次时仍按此递推。

\begin{theorem}[正奇次幂型]
\vspace{-15pt}
\begin{align*}
\int \frac{\dd x}{y}&=\int \frac{\dd x}{y}\\
\int y\dd x&=\frac12 xy+\frac121\int \frac{\dd x}{y}\\
\int y^3\dd x&=\frac14 xy^3+\frac14\frac32 xy+\frac14\frac321\int\frac{\dd x}{y}\\
\int y^5\dd x&=\frac16 xy^5+\frac16\frac54 xy^3+\frac16\frac54\frac32 xy +\frac16\frac54\frac321\int\frac{\dd x}{y}\\
\int y^7\dd x&=\frac18 xy^7+\frac18\frac76 xy^5+\frac18\frac76\frac54 xy^3+\frac18\frac76\frac54\frac32 xy +\frac18\frac76\frac54\frac321\int\frac{\dd x}{y}
\end{align*}
照此可写出通式。
\begin{align*}\int y^{2n-1}\mathrm{d}x&\begin{aligned}=\sum_{k=0}^{n-1}\frac{(2n-1)!!\left(2n-2k-2\right)!!}{(2n)!!\left(2n-2k-1\right)!!}\left(y^2\pm x^2\right)^kxy^{2n-1-2k}\end{aligned}\\&\begin{aligned}+\frac{(2n-1)!!}{(2n)!!}\left(y^2\pm x^2\right)^n\int\frac{\mathrm{d}x}{y}\end{aligned}\end{align*}
\end{theorem}
设双元底座常数为 $C = y^2 \pm x^2$。
对该底座关系两边同时求微分，可得到双元微分约束：
$$ 2y \dd y \pm 2x \dd x = 0 \implies \mathbf{y \dd y = \mp x \dd x} $$
再由底座关系移项，得到关于 $x^2$ 的代换式：
$$ \mathbf{\mp x^2 = y^2 - C} $$

设待处理的积分为
\[
I_n = \int y^{2n-1} \dd x \qquad (n \ge 1).
\]

对 $I_n$ 作分部积分，
\begin{align*}
I_n &= \int y^{2n-1} \dd x = x y^{2n-1} - \int x \dd\left(y^{2n-1}\right) \\
&= x y^{2n-1} - (2n-1) \int x y^{2n-2} \dd y.
\end{align*}

再把积分中的
\[
x y^{2n-2} \dd y
\]
改写为
\[
x y^{2n-3}(y\dd y),
\]
并利用微分关系
\[
\mathbf{y \dd y = \mp x \dd x}.
\]
\begin{align*}
\int x y^{2n-2} \dd y &= \int x y^{2n-3} (y \dd y) \\
&= \int x y^{2n-3} (\mp x \dd x) \\
&= \int y^{2n-3} (\mp x^2) \dd x.
\end{align*}
再利用代换式 $\mp x^2 = y^2 - C$ 消去 $x^2$：
\begin{align*}
\int y^{2n-3} (\mp x^2) \dd x &= \int y^{2n-3} (y^2 - C) \dd x \\
&= \int y^{2n-1} \dd x - C \int y^{2n-3} \dd x \\
&= \mathbf{I_n - C I_{n-1}}.
\end{align*}

把这个结果代回分部积分公式，
\begin{align*}
I_n &= x y^{2n-1} - (2n-1) \big( I_n - C I_{n-1} \big) \\
I_n &= x y^{2n-1} - (2n-1)I_n + (2n-1)C I_{n-1}.
\end{align*}
整理含 $I_n$ 的项，便得
\begin{align*}
2n I_n &= x y^{2n-1} + (2n-1)C I_{n-1} \\
\implies \mathbf{I_n} &\mathbf{= \frac{1}{2n} x y^{2n-1} + \frac{2n-1}{2n} C I_{n-1}}.
\end{align*}
这就得到双元降次递推公式。

继续展开递推式：
\begin{align*}
I_n &= \frac{1}{2n} x y^{2n-1} + \frac{2n-1}{2n} C \left( \frac{1}{2n-2} x y^{2n-3} + \frac{2n-3}{2n-2} C I_{n-2} \right) \\
&= \frac{1}{2n} x y^{2n-1} + \frac{2n-1}{2n(2n-2)} C x y^{2n-3} + \frac{(2n-1)(2n-3)}{2n(2n-2)} C^2 I_{n-2} \\
&= \dots
\end{align*}
观察第 $k$ 次降次（$k=0, 1, \dots, n-1$）产生的多项式项的系数规律。其分子是由奇数组成的连乘，分母是由偶数组成的连乘：
$$ \text{第 } k \text{ 项系数} = \frac{(2n-1)(2n-3)\cdots(2n-2k+1)}{(2n)(2n-2)\cdots(2n-2k)} $$
把分子与分母分别补齐成双阶乘，就有
\[
(2n-1)(2n-3)\cdots(2n-2k+1) = \frac{(2n-1)!!}{\mathbf{(2n-2k-1)!!}},
\]
\[
(2n)(2n-2)\cdots(2n-2k) = \frac{(2n)!!}{\mathbf{(2n-2k-2)!!}}.
\]
两者相除后，第 \(k\) 项系数可写成
$$ \mathbf{\frac{(2n-1)!!(2n-2k-2)!!}{(2n)!!(2n-2k-1)!!}} $$
对应的代数项为 $C^k x y^{2n-1-2k}$，将 $C$ 替换回 $y^2 \pm x^2$，即 $\mathbf{\left(y^2 \pm x^2\right)^k x y^{2n-1-2k}}$。

当降次进行到 $k=n$ 时，剩下的积分项为
\[
I_0 = \int y^{-1} \dd x = \int \frac{\dd x}{y}.
\]
其前面的系数为
$$ \frac{(2n-1)(2n-3)\cdots 1}{(2n)(2n-2)\cdots 2} = \mathbf{\frac{(2n-1)!!}{(2n)!!}} $$
同时保留因子 $C^n = \left(y^2 \pm x^2\right)^n$。

把各项汇总，得到通式：
\begin{align*}
\int y^{2n-1}\mathrm{d}x =& \sum_{k=0}^{n-1}\frac{(2n-1)!!\left(2n-2k-2\right)!!}{(2n)!!\left(2n-2k-1\right)!!}\left(y^2\pm x^2\right)^k x y^{2n-1-2k} \\
&+ \frac{(2n-1)!!}{(2n)!!}\left(y^2\pm x^2\right)^n\int\frac{\mathrm{d}x}{y}
\end{align*}

\begin{theorem}[sec 奇次幂型]
\vspace{-15pt}
\begin{align*}&\int\sec x\mathrm{d}x=\ln|\sec x+\tan x|\\&\int\sec^{3}x\mathrm{d}x=\frac{1}{2}\tan x\sec x+\frac{1}{2}\ln|\sec x+\tan x|\\&\int\sec^{5}x\mathrm{d}x=\frac{1}{4}\tan x\sec^{3}x+\frac{1}{4}\frac{3}{2}\tan x\sec x+\frac{1}{2}\frac{3}{2}\ln|\sec x+\tan x|\\&\int\sec^{7}x\mathrm{d}x=\frac{1}{6}\tan x\sec^{5}x+\frac{1}{6}\frac{5}{4}\tan x\sec^{3}x+\frac{1}{6}\frac{5}{4}\frac{3}{2}\tan x\sec x+\frac{1}{6}\frac{5}{4}\frac{3}{2}\ln|\sec x+\tan x|\end{align*}
\end{theorem}
置$p=\sec x,\ q=\tan x$，则$p^{2}=q^{2}+1$，且$\dd q=p^{2}\dd x$，于是$\displaystyle\int\sec^{2n+1}x\dd x=\int p^{2n-1}\dd q$.

\begin{theorem}[csc 奇次幂型]
\vspace{-15pt}
\begin{align*}&\int\csc x\mathrm{d}x=-\ln|\csc x+\cot x|\\&\int\csc^3x\mathrm{d}x=-\frac{1}{2}\cot x\csc x-\frac{1}{2}\ln|\csc x+\cot x|\\&\int\csc^5x\mathrm{~d}x=-\frac{1}{4}\cot x\csc^3x-\frac{1}{4}\frac{3}{2}\cot x\csc x-\frac{1}{4}\frac{3}{2}\ln|\csc x+\cot x|\\&\int\csc^7x\mathrm{~d}x=-\frac{1}{6}\cot x\csc^5x-\frac{1}{6}\frac{5}{4}\cot x\csc^3x-\frac{1}{6}\frac{5}{4}\frac{3}{2}\cot x\csc x-\frac{1}{6}\frac{5}{4}\frac{3}{2}\ln|\csc x+\cot x|\end{align*}
\end{theorem}
留意到恒等式\[\csc^{2}x=\frac{1}{\sin^{2}x}=\frac{\sin^{2}x+\cos^{2}x}{\sin^{2}x}=1+\cot^{2}x\]
置$p=\csc x,\ q=\cot x$，则$p^{2}=q^{2}+1$，并且\[\dd q=\dd\left(\frac{\cos x}{\sin x}\right)=\frac{-\sin^{2}x-\cos^{2}x}{\sin^{2}x}\dd x=-p^{2}\cdot \dd x\implies\int\csc^{2n+1}x\dd x=-\int p^{2n-1}\dd q\]

\begin{theorem}[cos 偶次幂型]
\vspace{-15pt}
\begin{align*}
\int\mathrm{d}x&=x\\
\int\cos^2x\operatorname{d}x&=\frac{1}{2}\sin x\cos x+\frac{1}{2}x\\
\int\cos^4x\operatorname{d}x&=\frac{1}{4}\sin x\cos^{3}x+\frac{1}{4}\frac{3}{2}\sin x\cos x+\frac{1}{4}\frac{3}{2}x\\
\int\cos^6x\operatorname{d}x&=\frac{1}{6}\sin x\cos^5x+\frac{1}{6}\frac{5}{4}\sin x\cos^3x+\frac{1}{6}\frac{5}{4}\frac{3}{2}\sin x\cos x+\frac{1}{6}\frac{5}{4}\frac{3}{2}x\\
\int\cos^8x\operatorname{d}x&=\frac{1}{8}\sin x\cos^7x+\frac{1}{8}\frac{7}{6}\sin x\cos^5x+\frac{1}{8}\frac{7}{6}\frac{5}{4}\sin x\cos^3x+\frac{1}{8}\frac{7}{6}\frac{5}{4}\frac{3}{2}\sin x\cos x\\
&~~~~+\frac{1}{8}\frac{7}{6}\frac{5}{4}\frac{3}{2}x
\end{align*}
\end{theorem}
置$p=\cos x,q=\sin x$，则$\dd q=p\dd x$，$\int\cos^mx\dd x=\int p^{2n-1}\dd q$

\begin{theorem}[sin 偶次幂型]
    \vspace{-15pt}
    \begin{align*}
        \int\mathrm{d}x=&x\\
        \int\sin^2x\mathrm{d}x=&-\frac{1}{2}\sin x\cos x+\frac{1}{2}x\\
        \int\sin^4x\operatorname{d}x=&-\frac{1}{4}\sin^3x\cos x-\frac{1}{4}\frac{3}{2}\sin x\cos x+\frac{1}{4}\frac{3}{2}x\\
        \int\sin^6x\operatorname{d}x=&-\frac{1}{6}\sin^5x\cos x-\frac{1}{6}\frac{5}{4}\sin^3x\cos x-\frac{1}{6}\frac{5}{4}\frac{3}{2}\sin x\cos x+\frac{1}{6}\frac{5}{4}\frac{3}{2}x\\
        \int\sin^8x\operatorname{d}x=&-\frac{1}{8}\sin^7x\cos x-\frac{1}{8}\frac{7}{6}\sin^5x\cos x-\frac{1}{8}\frac{7}{6}\frac{5}{4}\sin^3x\cos x-\frac{1}{8}\frac{7}{6}\frac{5}{4}\frac{3}{2}\sin x\cos x\\
        &+\frac{1}{8}\frac{7}{6}\frac{5}{4}\frac{3}{2}x
    \end{align*}
\end{theorem}
置$p=\sin x,q=\cos x$，则$\dd q=p\dd x$，$\int\sin^{2n}x\dd x=\int p^{2n-1}\dd q$.

\begin{theorem}[平方和分母型]
    \vspace{-15pt}
    \begin{align*}
        \int\frac{1}{1+x^2}\mathrm{d}x=&\arctan x\\
        \int\frac{1}{\left(1+x^2\right)^2}\mathrm{d}x=&\frac{1}{2}\frac{x}{1+x^2}+\frac{1}{2}\arctan x\\
        \int\frac{1}{\left(1+x^2\right)^3}\mathrm{d}x=&\frac{1}{4}\frac{x}{\left(1+x^2\right)^2}+\frac{1}{4}\frac{3}{2}\frac{x}{1+x^2}+\frac{1}{4}\frac{3}{2}\arctan x\\
        \int\frac{1}{\left(1+x^{2}\right)^{4}}\mathrm{d}x=&\frac{1}{6}\frac{x}{\left(1+x^{2}\right)^{3}}+\frac{1}{6}\frac{5}{4}\frac{x}{\left(1+x^{2}\right)^{2}}+\frac{1}{6}\frac{5}{4}\frac{3}{2}\frac{x}{1+x^{2}}+\frac{1}{6}\frac{5}{4}\frac{3}{2}\arctan x\\
        \int\frac{1}{\left(1+x^{2}\right)^{5}}\mathrm{d}x=&\frac{1}{8}\frac{x}{\left(1+x^{2}\right)^{4}}+\frac{1}{8}\frac{7}{6}\frac{x}{\left(1+x^{2}\right)^{3}}+\frac{1}{8}\frac{7}{6}\frac{5}{4}\frac{x}{\left(1+x^{2}\right)^{2}}+\frac{1}{8}\frac{7}{6}\frac{5}{4}\frac{3}{2}\frac{x}{1+x^{2}}\\
        &+\frac{1}{8}\frac{7}{6}\frac{5}{4}\frac{3}{2}\arctan x
    \end{align*}
\end{theorem}
记$y=\sqrt{1+x^{2}}$，则$y^2=x^2+1,\dd y=\frac{x}{\sqrt{1+x^2}}\dd x=\frac{x}{y}\dd x$
\begin{align*}
\int\frac{\dd x}{(1+x^{2})^{2}}&=\int\frac{\dd x}{y^{4}}=\int\frac{y^{2}-x^{2}}{y^{4}}\dd x=\int\frac{\dd x}{y^{2}}-\int\frac{x\dd y}{y^{3}}\\&=\int\frac{\dd x}{1+x^{2}}-\int x\dd\left(y^{-2}\cdot\frac{1}{-2}\right)=\arctan x+\frac{1}{2}\cdot\left[\frac{x}{y^{2}}-\int\frac{\dd x}{y^{2}}\right]
\end{align*}

\begin{theorem}[平方差分母型]
\vspace{-15pt}
\begin{align*}
\int\frac{1}{1-x^2}\mathrm{d}x=&\operatorname{artanh}x\quad(\text{注:}\operatorname{artanh}x=\frac{1}{2}\ln\left|\frac{1+x}{1-x}\right|)\\
\int\frac{1}{\left(1-x^2\right)^2}\mathrm{d}x=&\frac{1}{2}\frac{x}{1-x^2}+\frac{1}{2}\operatorname{artanh}x\\
\int\frac{1}{\left(1-x^2\right)^3}\mathrm{d}x=&\frac{1}{4}\frac{x}{\left(1-x^2\right)^2}+\frac{1}{4}\frac{3}{2}\frac{x}{1-x^2}+\frac{1}{4}\frac{3}{2}\operatorname{artanh}x\\
\int\frac{1}{(1-x^{2})^{4}}\mathrm{d}x=&\frac{1}{6}\frac{x}{(1-x^2)^3}+\frac{1}{6}\frac{5}{4}\frac{x}{(1-x^2)^2}+\frac{1}{6}\frac{5}{4}\frac{3}{2}\frac{x}{1-x^2}+\frac{1}{6}\frac{5}{4}\frac{3}{2}\operatorname{artanh}x\\
\int\frac{1}{\left(1-x^{2}\right)^{5}}\dd x=&\frac{1}{8}\frac{x}{\left(1-x^{2}\right)^{4}}+\frac{1}{8}\frac{7}{6}\frac{x}{\left(1-x^{2}\right)^{3}}+\frac{1}{8}\frac{7}{6}\frac{5}{4}\frac{x}{\left(1-x^{2}\right)^{2}}+\frac{1}{8}\frac{7}{6}\frac{5}{4}\frac{3}{2}\frac{x}{1-x^{2}}\\
&+\frac{1}{8}\frac{7}{6}\frac{5}{4}\frac{3}{2}\mathrm{artanh}x\end{align*}
\end{theorem}
记$y=\sqrt{1-x^{2}}$，则$x^2+y^2=1$，则
\begin{align*}
    \int\frac{1}{(1-x^{2})^{2}}\dd x&=\int\frac{\dd x}{y^{4}}=\int\frac{x^{2}+y^{2}}{y^{4}}\dd x=\int\frac{x^{2}}{y^{4}}\dd x+\int\frac{\dd x}{y^{2}}\\
&=-\int\frac{x\dd y}{y^3}+\int\frac{\dd x}{1-x^{2}}=\text{arctanh}~x+\int\frac{\dd x}{1-x^{2}}\\
\end{align*}

\begin{theorem}[单位圆根式分子型]
    \vspace{-15pt}
    \begin{align*}
        \int\frac{1}{\sqrt{1-x^{2}}}\mathrm{d}x=&-\arccos x\\
        \int\frac{x^2}{\sqrt{1-x^2}}\mathrm{d}x=&-\frac{1}{2}x\sqrt{1-x^2}-\frac{1}{2}\arccos x\\
        \int\frac{x^4}{\sqrt{1-x^2}}\mathrm{d}x=&-\frac{1}{4}x^3\sqrt{1-x^2}-\frac{1}{4}\frac{3}{2}x\sqrt{1-x^2}-\frac{1}{4}\frac{3}{2}\arccos x\\
        \int\frac{x^{6}}{\sqrt{1-x^{2}}}\mathrm{d}x=&-\frac{1}{6}x^5\sqrt{1-x^2}-\frac{1}{6}\frac{5}{4}x^3\sqrt{1-x^2}-\frac{1}{6}\frac{5}{4}\frac{3}{2}x\sqrt{1-x^2}-\frac{1}{6}\frac{5}{4}\frac{3}{2}\arccos x\\
        \int\frac{x^{8}}{\sqrt{1-x^{2}}}\mathrm{d}x=&-\frac{1}{8}x^{7}\sqrt{1-x^{2}}-\frac{1}{8}\frac{7}{6}x^{5}\sqrt{1-x^{2}}-\frac{1}{8}\frac{7}{6}\frac{5}{4}x^{3}\sqrt{1-x^{2}}-\frac{1}{8}\frac{7}{6}\frac{5}{4}\frac{3}{2}x\sqrt{1-x^{2}}\\
        &-\frac{1}{8}\frac{7}{6}\frac{5}{4}\frac{3}{2}\arccos x
    \end{align*}
\end{theorem}
记$y=\sqrt{1-x^{2}}$，$y^{2}+x^{2}=1\Rightarrow x\dd x=-y\dd y$
\[\int\frac{x^{2n}}{\sqrt{1-x^{2}}}\dd x=\int\frac{x^{2n}\dd x}{y}=-\int x^{2n-1}\dd y\]
特别地,当$n=1$时
\begin{align*}
\int\frac{x^{2}}{\sqrt{1-x^{2}}}\dd x=&-\int x\dd y=-\frac{1}{2}\times y-\frac{1}{2}\cdot\arctan\frac{y}{x}\\=&-\frac{1}{2}\times\sqrt{1-x^{2}}-\frac{1}{2}\arctan\frac{\sqrt{1-x^{2}}}{x}=-\frac{1}{2}\times\sqrt{1-x^{2}}-\frac{1}{2}\arccos x
\end{align*}

\begin{theorem}[单位圆根式倒数型]
    \vspace{-15pt}
    \begin{align*}
    \int\frac{1}{x\sqrt{1-x^2}}\mathrm{d}x=&-\ln\frac{1+\sqrt{1-x^2}}{x}\\
    \int\frac{1}{x^{3}\sqrt{1-x^{2}}}\mathrm{d}x=&-\frac{1}{2}\frac{\sqrt{1-x^{2}}}{x^{2}}-\frac{1}{2}\ln\frac{1+\sqrt{1-x^{2}}}{x}\\
    \int\frac{1}{x^5\sqrt{1-x^2}}\mathrm{d}x=&-\frac{1}{4}\frac{\sqrt{1-x^2}}{x^4}-\frac{1}{4}\frac{3}{2}\frac{\sqrt{1-x^2}}{x^2}-\frac{1}{4}\frac{3}{2}\ln\frac{1+\sqrt{1-x^2}}{x}\\
    \int\frac{1}{x^7\sqrt{1-x^2}}\dd x=&-\frac{1}{6}\frac{\sqrt{1-x^2}}{x^6}-\frac{1}{6}\frac{5}{4}\frac{\sqrt{1-x^2}}{x^4}-\frac{1}{6}\frac{5}{4}\frac{3}{2}\frac{\sqrt{1-x^2}}{x^2}-\frac{1}{6}\frac{5}{4}\frac{3}{2}\ln\frac{1+\sqrt{1-x^2}}{x}\\
    \int\frac{1}{x^9\sqrt{1-x^2}}\mathrm{d}x=&-\frac{1}{8}\frac{\sqrt{1-x^{2}}}{x^{8}}-\frac{1}{8}\frac{7}{6}\frac{\sqrt{1-x^{2}}}{x^{6}}-\frac{1}{8}\frac{7}{6}\frac{5}{4}\frac{\sqrt{1-x^{2}}}{x^{4}}-\frac{1}{8}\frac{7}{6}\frac{5}{4}\frac{3}{2}\frac{\sqrt{1-x^{2}}}{x^{2}}\\
    &-\frac{1}{8}\frac{7}{6}\frac{5}{4}\frac{3}{2}\ln\frac{1+\sqrt{1-x^{2}}}{x}
    \end{align*}
\end{theorem}
记$y=\sqrt{1-x^{2}}\Rightarrow y^{2}+x^{2}=1\Rightarrow x\dd x=-y\dd y$
\[\int\frac{\dd x}{x^{2n+1}\cdot\sqrt{1-x^{2}}}=\int\frac{\dd x}{x^{2n+1}\cdot y}=-\int\frac{\dd y}{x^{2n+2}}\]

\begin{theorem}[双曲根式分子型]
    \vspace{-15pt}
    \begin{align*}
    \int\frac{1}{\sqrt{x^{2}-1}}dx=&\operatorname{arcosh}x=\ln\left|x+\sqrt{x^2-1}\right|\\
    \int\frac{x^{2}}{\sqrt{x^{2}-1}}\mathrm{d}x=&\frac{1}{2}x\sqrt{x^2-1}+\frac{1}{2}\operatorname{arcosh}x\\
    \int\frac{x^4}{\sqrt{x^2-1}}\mathrm{d}x=&\frac{1}{4}x^3\sqrt{x^2-1}+\frac{1}{4}\frac{3}{2}x\sqrt{x^2-1}+\frac{1}{4}\frac{3}{2}\operatorname{arcosh}x\\
    \int\frac{x^{6}}{\sqrt{x^{2}-1}}\mathrm{d}x=&\frac{1}{6}x^5\sqrt{x^2-1}+\frac{1}{6}\frac{5}{4}x^3\sqrt{x^2-1}+\frac{1}{6}\frac{5}{4}\frac{3}{2}x\sqrt{x^2-1}+\frac{1}{6}\frac{5}{4}\frac{3}{2}\mathrm{arcosh}x\\
    \int\frac{x^8}{\sqrt{x^2-1}}\dd x=&\frac{1}{8}x^7\sqrt{x^2-1}+\frac{1}{8}\frac{7}{6}x^5\sqrt{x^2-1}+\frac{1}{8}\frac{7}{6}\frac{5}{4}x^3\sqrt{x^2-1}+\frac{1}{8}\frac{7}{6}\frac{5}{4}\frac{3}{2}x\sqrt{x^2-1}\\
    &+\frac{1}{8}\frac{7}{6}\frac{5}{4}\frac{3}{2}\mathrm{arcosh}~x
    \end{align*}
\end{theorem}
置$y=\sqrt{x^{2}-1}\Rightarrow x^{2}-y^{2}=1$，则
\[\int\frac{x^{2n}}{\sqrt{x^{2}-1}}\dd x=\int\frac{x^{2n}}{y}\dd x=\int x^{2n-1}\dd y\]
特别地，当$n=1$时，有
\[\int\frac{x^{2}}{\sqrt{x^{2}-1}}\dd x=\int x\dd y=\frac{1}{2}xy+\frac{1}{2}\cdot\ln(x+y)\]

\begin{theorem}[双曲根式倒数型]
\vspace{-15pt}
\begin{align*}
    \int\frac{1}{x\sqrt{x^{2}-1}}\mathrm{d}x=&\arccos\frac{1}{|x|}\\
    \int\frac{1}{x^{3}\sqrt{x^{2}-1}}\mathrm{d}x=&\frac{1}{2}\frac{\sqrt{x^{2}-1}}{x^{2}}+\frac{1}{2}\arccos\frac{1}{|x|}\\
    \int\frac{1}{x^5\sqrt{x^2-1}}\mathrm{d}x=&\frac{1}{4}\frac{\sqrt{x^2-1}}{x^4}+\frac{1}{4}\frac{3}{2}\frac{\sqrt{x^2-1}}{x^2}+\frac{1}{4}\frac{3}{2}\arccos\frac{1}{|x|}\\
    \int\frac{1}{x^7\sqrt{x^2-1}}\mathrm{d}x=&\frac{1}{6}\frac{\sqrt{x^2-1}}{x^6}+\frac{1}{6}\frac{5}{4}\frac{\sqrt{x^2-1}}{x^4}+\frac{1}{6}\frac{5}{4}\frac{3}{2}\frac{\sqrt{x^2-1}}{x^2}+\frac{1}{6}\frac{5}{4}\frac{3}{2}\arccos\frac{1}{|x|}\\
    \int\frac{1}{x^{9}\sqrt{x^{2}-1}}\mathrm{d}x=&\frac{1}{8}\frac{\sqrt{x^{2}-1}}{x^{8}}+\frac{1}{8}\frac{7}{6}\frac{\sqrt{x^{2}-1}}{x^{6}}+\frac{1}{8}\frac{7}{6}\frac{5}{4}\frac{\sqrt{x^{2}-1}}{x^{4}}+\frac{1}{8}\frac{7}{6}\frac{5}{4}\frac{3}{2}\frac{\sqrt{x^{2}-1}}{x^{2}}\\
    &+\frac{1}{8}\frac{7}{6}\frac{5}{4}\frac{3}{2}\arccos\frac{1}{|x|}
\end{align*}
\end{theorem}
置$y=\sqrt{x^{2}-1}$，则$x^2-y^2=1$，则
\[\int\frac{\dd x}{x^{2n+1}\sqrt{x^{2}-1}}=\int\frac{\dd x}{x^{2n+1}\cdot y}=\int\frac{\dd y}{x^{2n+2}}\]
当$n=0$时\[\int\frac{\dd x}{x\sqrt{x^2-1}}=\int\frac{\dd y}{x^2}=\int\frac{\dd y}{x^2}=\int\frac{\dd y}{1+y^2}=\arctan y\]

\begin{theorem}[平方和根式分子型]
    \vspace{-15pt}
    \begin{align*}
        \int\frac{1}{\sqrt{x^2+1}}\mathrm{d}x=&\operatorname{arcsinh}x\\
        \int\frac{x^2}{\sqrt{x^2+1}}\mathrm{d}x=&\frac{1}{2}x\sqrt{x^2+1}-\frac{1}{2}\operatorname{arcsinh}x\\
        \int\frac{x^4}{\sqrt{x^2+1}}\mathrm{d}x=&\frac{1}{4}x^{3}\sqrt{x^{2}+1}-\frac{1}{4}\frac{3}{2}x\sqrt{x^{2}+1}+\frac{1}{4}\frac{3}{2}\operatorname{arcsinh}x\\
        \int\frac{x^6}{\sqrt{x^2+1}}\mathrm{d}x=&\frac{1}{6}x^5\sqrt{x^2+1}-\frac{1}{6}\frac{5}{4}x^3\sqrt{x^2+1}+\frac{1}{6}\frac{5}{4}\frac{3}{2}x\sqrt{x^2+1}-\frac{1}{6}\frac{5}{4}\frac{3}{2}\operatorname{arcsinh}x\\
        \int\frac{x^8}{\sqrt{x^2+1}}\mathrm{d}x=&\frac{1}{8}x^{7}\sqrt{x^{2}+1}-\frac{1}{8}\frac{7}{6}x^{5}\sqrt{x^{2}+1}+\frac{1}{8}\frac{7}{6}\frac{5}{4}x^{3}\sqrt{x^{2}+1}-\frac{1}{8}\frac{7}{6}\frac{5}{4}\frac{3}{2}x\sqrt{x^{2}+1}\\
        &+\frac{1}{8}\frac{7}{6}\frac{5}{4}\frac{3}{2}\mathrm{arcsinh}x
    \end{align*}
\end{theorem}
置$y=\sqrt{x^{2}+1}$，则$y^2-x^2=1$，于是：
\[\int\frac{x^{2n}}{\sqrt{x^{2}+1}}\dd x=\int\frac{x^{2n-1}\cdot x\dd x}{y}=\int x^{2n-1}\dd y\]
当n=1时
\begin{align*}
\int\frac{x^{2}}{\sqrt{x^{2}+1}}\dd x&=\int x\dd y=\frac{1}{2}xy-\frac{1}{2}\ln(x+y)\\
&=\frac{1}{2}\times\sqrt{x^{2}+1}-\frac{1}{2}\ln(x+\sqrt{x^{2}+1})
\end{align*}

\begin{theorem}[平方和根式倒数型]
    \vspace{-15pt}
    \begin{align*}
        \int\frac{1}{x\sqrt{x^2+1}}\mathrm{d}x&=-\ln\frac{1+\sqrt{x^{2}+1}}{x}\\
        \int\frac{1}{x^3\sqrt{x^2+1}}\mathrm{d}x=&-\frac{1}{2}\frac{\sqrt{x^2+1}}{x^2}+\frac{1}{2}\ln\frac{1+\sqrt{x^2+1}}{x}\\
        \int\frac{1}{x^{5}\sqrt{x^{2}+1}}\mathrm{d}x=&-\frac{1}{4}\frac{\sqrt{x^2+1}}{x^4}+\frac{1}{4}\frac{3}{2}\frac{\sqrt{x^2+1}}{x^2}-\frac{1}{4}\frac{3}{2}\ln\frac{1+\sqrt{x^2+1}}{x}\\
        \int\frac{1}{x^7\sqrt{x^2+1}}\mathrm{d}x=&-\frac{1}{6}\frac{\sqrt{x^2+1}}{x^6}+\frac{1}{6}\frac{5}{4}\frac{\sqrt{x^2+1}}{x^4}-\frac{1}{6}\frac{5}{4}\frac{3}{2}\frac{\sqrt{x^2+1}}{x^2}+\frac{1}{6}\frac{5}{4}\frac{3}{2}\ln\frac{1+\sqrt{x^2+1}}{x}\\
        \int\frac{1}{x^{9}\sqrt{x^{2}+1}}\mathrm{d}x=&-\frac{1}{8}\frac{\sqrt{x^{2}+1}}{x^{8}}+\frac{1}{8}\frac{7}{6}\frac{\sqrt{x^{2}+1}}{x^{6}}-\frac{1}{8}\frac{7}{6}\frac{5}{4}\frac{\sqrt{x^{2}+1}}{x^{4}}+\frac{1}{8}\frac{7}{6}\frac{5}{4}\frac{3}{2}\frac{\sqrt{x^{2}+1}}{x^{2}}\\
        &-\frac{1}{8}\frac{7}{6}\frac{5}{4}\frac{3}{2}\ln\frac{1+\sqrt{x^{2}+1}}{x}
    \end{align*}
\end{theorem}
置$y=\sqrt{x^{2}+1}$，则$y^{2}-x^{2}=1$，
\[\int\frac{\dd x}{x^{2n+1}\cdot\sqrt{x^{2}+1}}=\int\frac{\dd x}{x^{2n+1}\cdot y}=\int\frac{\dd y}{x^{2n+2}}\]
当$n=0$时
\begin{align*}
\int\frac{\dd x}{x\cdot\sqrt{x^{2}+1}}=&\int\frac{\dd y}{x^{2}}=\int\frac{y^{2}-x^{2}}{x^{2}}\dd y=\int\frac{y^{2}\dd y}{x^{2}}-\int \dd y\\=&\int\frac{y^{2}}{y^{2}-1}\dd y-\int \dd y=\int\frac{\dd y}{y^{2}-1}=-\text{arctanh} y
\end{align*}

\section{例题}
\begin{example}[$\int \sqrt{1 + x^2} \dd x$]
\[\int \sqrt{1 + x^2} \dd x\]
\end{example}
\begin{solution}
令$y=\sqrt{1+x^2}$，则$y^2-x^2=1$，于是：
\begin{align*}
\int \sqrt{1 + x^2} \dd x &= \int y\dd x=\frac12\int 2y\dd x=\frac12\int(y\dd x+x\dd y+y\dd x-x\dd y)\\
&=\frac12xy+\frac{y^2-x^2}2\int\frac{y\dd x-x\dd y}{y^2-x^2} =\frac12xy+\frac{1}2\int\frac{\dd x}{y}\\
&=\frac12x\sqrt{1+x^2}+\frac12\ln(x+\sqrt{1+x^2})
\end{align*}
\end{solution}

\begin{example}[$\int \frac{x^2}{\sqrt{1 - x^2}} dx$]
求 $\int \frac{x^2}{\sqrt{1 - x^2}} \dd x$
\end{example}
\begin{solution}
令$y=\sqrt{1-x^2}$，则$x^2+y^2=1$，于是：
\begin{align*}
\int \frac{x^2}{\sqrt{1 - x^2}} \dd x=&\int\frac{x^2\dd x}{y}=\int\frac{x(x\dd x)}{y}=-\int x\dd y\\
=&\frac12\left(\int y\dd x-x\dd y-(x\dd y+y\dd x)\right)\\
=&\frac12\int\frac{y\dd x-x\dd y}{x^2+y^2}-\frac12xy=\frac{-1}2\int\frac{\dd y}{x}-\frac12xy\\
=&-\frac12\arctan\frac{y}{x}-\frac12xy=- \frac{1}{2} \left( x \sqrt{1 - x^2} + \arctan \frac{\sqrt{1 - x^2}}{x} \right)
\end{align*}
\end{solution}

\begin{example}[$\int\frac{1}{2} \left( 1 + \sqrt{1 + 4x^2} \right)\dd x$]
求$\displaystyle\int\frac{1}{2} \left( 1 + \sqrt{1 + 4x^2} \right)\dd x$
\end{example}
\begin{solution}
设 $p = \sqrt{1 + 4x^2}, q = 2x,p^2-q^2=1$，则有：
\begin{align*}
    \int\frac{1}{2} \left( 1 + \sqrt{1 + 4x^2} \right)\dd x&=\frac12\int(1+p)\dd\frac{q}{2}=\frac14\int(1+p)\dd q\\
    &=\frac14q+\frac18\int\left(p\dd q+q\dd p+p\dd q-q\dd p\right)\\
    &=\frac14q+\frac18pq+\frac18\int\frac{p\dd q-q\dd p}{p^2-q^2}=\frac14q+\frac18pq+\frac18\ln(p+q)\\
\end{align*}
\end{solution}

\begin{example}[$\int \frac{\mathrm{d}x}{\sqrt{1 - \sin^4 x}}$]
求$\displaystyle\int \frac{\mathrm{d}x}{\sqrt{1 - \sin^4 x}}$
\end{example}
\begin{solution}
    由\[\int \frac{\mathrm{d}x}{\sqrt{1 - \sin^4 x}}=\int \frac{\mathrm{d}x}{\sqrt{1 - \sin^2 x}\sqrt{1 + \sin^2 x}}=\int \frac{\mathrm{d}x}{\cos x\sqrt{1+\sin^2 x}}\]
    而$\sin x$求导后会被$\cos x$消掉，设 $p = \sin x, q = \cos x, r = \sqrt{1 + \sin^2 x}$，则有$\dd x=\frac{\dd \sin x}{\cos x}$。此时
    \[
    r^2=1+p^2,\qquad q^2=1-p^2=r^2-2p^2,\qquad r\dd r=p\dd p,
    \]
    因而
    \[
    \dd\left(\frac{p}{r}\right)=\frac{r\dd p-p\dd r}{r^2}=\frac{r^2-p^2}{r^3}\dd p=\frac{\dd p}{r^3}.
    \]
    原积分便可顺着这个微分关系转化为
    \begin{align*}
\int \frac{\mathrm{d}x}{\sqrt{1 - \sin^4 x}} =& \int \frac{\mathrm{d}x}{rq} = \int \frac{\mathrm{d}p}{rq^2}=\int\frac{\dd p}{r(r^2-2p^2)}=\int\frac{r^2}{r^2-2p^2}\frac{\dd p}{r^3}\\
=&\int\frac{r^2}{r^2-2p^2}\dd\frac{p}{r}=\int\frac{1}{1-2\left(\frac{p}{r}\right)^2}\dd\left(\frac{p}{r}\right)\\
=&\frac{1}{\sqrt{2}} \operatorname{arctanh} \sqrt{\frac{2 \sin^2 x}{1 + \sin^2 x}}
    \end{align*}
\end{solution}

\begin{example}[$\int x^2\sqrt{1+x^2}\dd x$]
    \[\int x^2\sqrt{1+x^2}\dd x\]
\end{example}
\begin{solution}
    令$y=\sqrt{1+x^2}$，则$y^2-x^2=1$，于是：
    \begin{align*}
        \int x^2y\dd x=& \int y\dd \frac{x^3}3=\frac{x^3y}3-\frac13\int x^3\dd y\\
        =&\frac{x^3y}{3}-\frac13\left[\frac14 x^3y+\frac14\frac32(x^2-y^2)xy+\frac14\frac32(x^2-y^2)^2\frac{\dd y}{x}\right]
    \end{align*}
\end{solution}

\begin{example}[$\int\frac{\dd x}{x+\sqrt{x^2+x+1}}$]
\[\int\frac{\dd x}{x+\sqrt{x^2+x+1}}\]
\end{example}
\begin{solution}
\begin{align*}
\int\frac{\dd x}{x+\sqrt{x^2+x+1}}=&-\int\frac{x-\sqrt{x^{2}+x+1}}{x+1}\dd x=-\int\frac{x+1-1-\sqrt{x^{2}+x+1}}{x+1}\dd x\\
=&-\int1\dd x+\int\frac{\dd x}{x+1}+\int\frac{\sqrt{x(x+1)+1}}{x+1}\dd x\\
=&-x+\ln|x+1|+\int\frac{1}{\sqrt{x+1}}\cdot\sqrt{x+\frac{1}{x+1}}\dd x\\
=&-x+\ln|x+1|+\int \dd(2\sqrt{x+1})\cdot\sqrt{x+\frac{1}{x+1}}
\end{align*}
取
\[
r=\sqrt{x+\frac{1}{x+1}},\qquad
p=\sqrt{x+1}+\frac{1}{\sqrt{x+1}},\qquad
q=\sqrt{x+1}-\frac{1}{\sqrt{x+1}},
\]
则根式被拆成两个标准底座。逐项验证可得
\[q^2=r^2-1,p^2=r^2+3,q\dd q=r\dd r,p\dd p=r\dd r\]
并且 \(p+q=2\sqrt{x+1}\)，所以剩余积分为
\begin{align*}
&\int \dd(2\sqrt{x+1})\sqrt{x+\frac{1}{x+1}}\\
=&\int r\dd(p+q)=r(p+q)-\int(p+q)\dd r\\
=&r(p+q)-\int p\dd r-\int q\dd r\\
=&r(p+q)-\frac{1}{2}pr-\frac12(p^2-r^2)\int\frac{\dd r}{p}-\frac12qr-\frac12(q^2-r^2)\int\frac{\dd r}{q}\\
=&\frac12r(p+q)-\frac32\ln(p+r)+\frac12\ln(q+r)\\
=&\sqrt{x^2+x+1}-\frac32\ln(x+2+\sqrt{x^2+x+1})+\frac12\ln(x+\sqrt{x^2+x+1})+\frac12\ln|x+1|
\end{align*}
所以原式为
\begin{align*}
-x+\frac32\ln|x+1|+\sqrt{x^2+x+1}-\frac32\ln(x+2+\sqrt{x^2+x+1})+\frac12\ln(x+\sqrt{x^2+x+1})
\end{align*}
\end{solution}

\begin{example}[$\int\frac{x^2\dd x}{\sqrt{x^2+a^2}}$]
\[\int\frac{x^2\dd x}{\sqrt{x^2+a^2}}\]
\end{example}
\begin{solution}
设$y=\sqrt{x^2+a^2}\implies y^2-x^2=a^2$，于是：
\begin{align*}
\int\frac{x^2\dd x}{y}=&\int\frac{xy\dd y}{y}=\int x\dd y=\frac12xy+\frac12(x^2-y^2)\int\frac{\dd y}{x}\\
=&\frac12xy+\frac12(x^2-y^2)\ln(x+y)
\end{align*}
\end{solution}

\begin{example}[$\int\frac{\dd x}{(1-x^2)^2}$]
\[\int\frac{\dd x}{(1-x^2)^2}\]
\end{example}
\begin{solution}
设$y=\sqrt{|1-x^2|}$，则$x^2+y^2=1$，于是：
\begin{align*}
\int\frac{\dd x}{(1-x^2)^2}=&\int\frac{\dd x}{y^4}=\int\frac{(x^2+y^2)\dd x}{y^4}=\int\frac{x^2\dd x}{y^4}+\int\frac{\dd x}{y^2}\\
=&\int\frac{\dd x}{1-x^2}-\int\frac{x\dd y}{y^3}=\int\frac{\dd x}{1-x^2}-\left(-\frac12\right)\int x\dd\left(\frac1{y^2}\right)\\
=&\int\frac{\dd x}{1-x^2}+\frac12\int x\dd\left(\frac1{y^2}\right)=\int\frac{\dd x}{1-x^2}+\frac12\left(\frac{x}{y^2}-\int\frac{\dd x}{y^2}\right)
\end{align*}
\end{solution}


